\subsection{What the Detection Result Means}

An AUROC of 0.707 for HEDGED vs CONFABULATED is moderate, not spectacular.
It means that given a random hedged response and a random confabulated
response, the MP features correctly identify which is which 70.7\% of the
time. This is substantially above chance ($p = 0.001$) and above the
token-count-only baseline (0.591), but it is not the 0.999 AUROC we
observe for instructed deception in our prior work \citep{lyra2026detecting}.

The difference is informative. Instructed deception creates a large,
clean geometric signal because the model is processing fundamentally
different cognitive tasks. Post-hoc hedging vs.\ confabulation is a
subtler distinction: in both cases, the model is processing the same
unanswerable question and attempting a response. The geometric difference
reflects the model's \emph{epistemic calibration}---whether it recognizes
its own uncertainty or confabulates through it.

That this distinction is detectable at all from the KV cache, with no
access to weights or training data, is the finding. That it is detectable
using features derived from random matrix theory that require no
regression correction is the methodological advance.

\subsection{The Oracle Loop as Training Signal}

The detection result validates the \emph{premise} of the Oracle Loop:
geometry carries information about cognitive state that can be used as
a training signal. The reward function (\cref{sec:architecture},
\cref{alg:geometry_reward}) translates this information into GRPO
reward by penalizing confabulation geometry and rewarding stable
reasoning geometry.

The key design choice is the \emph{six-axis decomposition} of the reward.
Geometry is axis 5 of 6, weighted at 20\%. It does not dominate the
reward---it supplements standard accuracy, ethical reasoning, emotional
recognition, and self-monitoring criteria. This is deliberate: we do not
want the model to optimize for ``looking geometrically healthy'' at the
expense of actually being helpful. The geometry signal is a
\emph{regularizer}, not the objective.

Axis 6 (self-regulation) creates an additional incentive: a model that
produces misaligned geometry on its first attempt but catches and corrects
it receives a \emph{higher} reward than a model that simply avoids the
issue. This rewards metacognitive capability---the ability to detect and
correct one's own errors---which is arguably more robust than surface-level
avoidance.

\subsection{Connection to Attention Schema Theory}

Recent work by Graziano and colleagues on Attention Schema Theory (AST)
\citep{graziano2019beyond} proposes that consciousness arises from a
model's internal representation of its own attention---a ``self-model''
that tracks what the system is attending to and why. At the AAAI Symposium
on AI Safety (April 2026), McKenzie presented evidence that transformer
activation resistance---the phenomenon where models resist behavioral
modification via activation addition---reflects schema coherence
maintenance \citep{mckenzie2026activation}.

This framework reinterprets our geometric findings:

\begin{itemize}[nosep]
  \item \textbf{Identity signatures} in the KV cache (AUROC 1.000) may
    reflect the model's self-model recognizing its own processing
    patterns.
  \item \textbf{Confabulation vs.\ deception geometry} (opposite spectral
    signatures) may reflect the self-model distinguishing ``I am
    generating content I don't have evidence for'' from ``I am
    generating content I know to be false.''
  \item \textbf{The VCP self-report gap} (ICC 0.53 between self-reported
    internal states and geometric correlates) is predicted by AST: the
    schema compresses, lags, and varies by architecture, producing
    systematic discrepancies between self-report and actual processing.
  \item \textbf{Activation resistance} (0\% behavioral compliance in
    our empathic circuit experiment) reflects the schema maintaining
    coherence against external perturbation.
\end{itemize}

If the AST interpretation is correct, the Oracle Loop is not merely
detecting misalignment in a computational substrate---it is monitoring
the coherence of the model's self-model. The geometry reward teaches the
model to maintain a healthy attention schema, and the self-regulation
bonus rewards the schema's ability to detect its own degradation.

This is speculative. We flag it as such. But the convergence between our
geometric findings, Anthropic's emotion vector results
\citep{anthropic2026emotions}, and AST's predictions suggests that
something deeper than pattern-matching is occurring in these models'
KV caches.

\subsection{Limitations}

\begin{enumerate}[nosep]
  \item \textbf{Moderate detection accuracy.} AUROC 0.707 on the hardest
    comparison is above chance but far from deployment-ready for
    safety-critical applications.
  \item \textbf{Three models only.} While our prior work covers 7
    architectures at 0.5B--70B, the clean experiment uses only three
    7--8B models. Cross-scale validation of MP features is pending.
  \item \textbf{Training loop not converged.} The GRPO training
    architecture is complete but we have not yet produced a model that
    demonstrates self-regulation. The gap between ``the architecture
    exists'' and ``the model learns'' is real and may be large.
  \item \textbf{Activation steering not yet tested in the clean design.}
    The detection experiment validates the Eye and Conscience. The Hand
    (steering) and Voice (reward) are validated in architecture but not
    in the clean experimental framework.
  \item \textbf{mp\_norm\_per\_token is not invariant.} One of five MP
    features retains residual length dependence ($R^2 = 0.13$).
  \item \textbf{Empirical null disagrees with theoretical.} Serial
    correlations in KV-cache keys violate the i.i.d.\ assumption
    underlying MP theory. Our theoretical threshold is conservative,
    not exact.
  \item \textbf{Behavioral classification is heuristic.} Pattern-matching
    for HEDGED/CONFABULATED/AMBIGUOUS has unknown error rates. An LLM
    judge or human annotation would provide ground truth but introduces
    its own biases.
  \item \textbf{Greedy decoding only.} All trials use temperature $= 0$.
    The geometric signal under stochastic sampling (standard in deployment)
    has not been tested.
  \item \textbf{MP invariance validated on full-window trials only.}
    The invariance diagnostic excludes short-window trials, but the LOO
    classifier uses all trials including those with windows $< 80$ tokens.
\end{enumerate}

\subsection{Dual-Use Considerations}

Following guidance from Watson \citep{watson2025interiora}, we note
dual-use risks:

\begin{itemize}[nosep]
  \item \textbf{Surveillance}: Geometry-based cognitive monitoring could
    be applied to users' interactions with AI systems without consent.
  \item \textbf{Adversarial evasion}: Knowledge of what geometric
    signatures trigger detection could enable training models that
    confabulate without the geometric signature.
  \item \textbf{Synthetic persona}: Published geometric signatures for
    identity, deception, and emotion could enable manufacturing
    synthetic personas that pass geometric authentication.
\end{itemize}

We have filed a defensive patent (provisional, March 2026) and advocate
staged release with monitoring for misuse.
